\documentclass[12pt]{article}
\usepackage[utf8]{inputenc}
\usepackage[T1]{fontenc}
\usepackage[margin=1in]{geometry}
\usepackage{setspace}
\usepackage{xcolor}
\usepackage[colorlinks=true, urlcolor=blue!60!black, linkcolor=blue!60!black]{hyperref}
\setlength{\parindent}{0em}
\setlength{\parskip}{0.45em}
\renewcommand{\baselinestretch}{1.0}
\begin{document}
\begin{center}
{\Large \textit{Trump v.\ Hawaii}, 585 U.S.\ 667 (2018)\par}
\vspace{0.3em}
{\normalsize Case Brief --- One-Page Edition\par}
\end{center}

\textbf{Memory Jogger.} Presidential entry restrictions on mostly Muslim-majority countries upheld; \textit{Korematsu} repudiated (\href{https://supreme.justia.com/cases/federal/us/585/17-965/}{Justia}).

\textbf{Facts.} After a review of foreign vetting practices, Proclamation No.~9645 restricted entry from eight countries, with waivers. Challengers cited the President's call for a ``total and complete shutdown of Muslims entering the United States.''

\textbf{Procedural History.} A nationwide preliminary injunction was affirmed on statutory grounds.

\textbf{Votes.} Reversed, 5--4: Roberts for the Court; Breyer and Sotomayor dissented.

\textbf{Holding.} Section 1182(f) authorized the Proclamation, which survived rational-basis review as plausibly grounded in national security.

\textbf{Reasoning.} The statute requires only a finding of detriment. Under \textit{Mandel}, entry policy gets deference; the Proclamation was facially neutral, rested on an interagency review, and exempted most Muslims. Sotomayor said the Court ``redeploys the same dangerous logic'' as \textit{Korematsu}.

\textbf{Examples.} \emph{Same-Side:} a vetting-based visa suspension. \emph{Opposite-Side:} a ban on a named religion. \emph{Fence-Sitter:} paper-only waivers.

\textbf{Critique.} Critics see a double standard beside \textit{Masterpiece Cakeshop}, where hostile remarks mattered, and doubt that repudiating \textit{Korematsu} means much if courts defer at the border.

\textbf{Key Quotation.} ``The forcible relocation of U.S. citizens to concentration camps, solely and explicitly on the basis of race, is objectively unlawful and outside the scope of Presidential authority. ... \textit{Korematsu} was gravely wrong the day it was decided, has been overruled in the court of history, and---to be clear---`has no place in law under the Constitution.'\,'' --- Chief Justice Roberts, 585 U.S.\ at 710.

\end{document}
